\section{Exclusion of every rank-three linear and unitary core}
\label{sec:a2-cores}

\paragraph{Status and exact scope.}
This is a new proof from complete published character tables, their
corrected official mathematical implementation, and the published
classification of bicyclic extensions of $\operatorname{PSL}_3(4)$.
The full proof has passed the root reading and an independent whole-proof
audit by \texttt{/root/g34\_eightfolds}.  That audit independently
reconstructed every generic inequality from the original frozen expressions
and every finite spectrum from the complete retained tables.  The separate
L34 normalizer proof also passed a whole-proof audit by
\texttt{/root/imprimitive}.  The precise source and computation scopes
are recorded in their reports and the check-event register.
All field sizes $q\ge3$ are covered.  The simple group
$\operatorname{PSL}_3(2)\simeq\operatorname{PSL}_2(7)$ is expressly
outside the statement: its degree-three representation occurs in the
already successful split-sixfold construction.  The group
$\operatorname{PSU}_3(2)$ is not simple.

For a finite-order operator on a complex vector space of dimension $a$,
put $\nu(T)=a-\max_\lambda\dim\ker(T-\lambda)$.
Scalar multiplication does not change $\nu$.

\subsection*{D-A2-01. A uniform bound on every ordinary character}

\begin{lemma}[D-A2-01]
\label{D-A2-01}
Let $U=\operatorname{SL}_3(q)$, $q\ge5$, or
$U=\operatorname{SU}_3(q)$, $q\ge7$.  For every nontrivial ordinary
irreducible character $\chi$ and every $u\notin Z(U)$,
\[
                 |\chi(u)|\le\chi(1)/4.                 \tag{A2.1}
\]
The minimum nontrivial degree is at least $30$ in the linear range
and at least $42$ in the unitary range.
\end{lemma}

\begin{proof}
Use the complete Simpson--Frame tables, Tables~1a, 1b and~2
\cite{A2:SF73}, with the unitary corrections of Geck,
\S3 \cite{A2:Geck90}, as implemented in the four complete official
tables of \cite{G2T:GCT}.  We freeze the implementation at commit
\texttt{87231738271ed548f62e7f1c98ff2bbfc1ed30d3}.
The inventory is
\[
\begin{array}{l|l|r|r|r|r}
\text{table}&\text{field condition}&q_0&\text{row types}&
 \text{class types}&\min\chi(1)\\\hline
\texttt{SL3.1}&q\equiv1\pmod3&5&17&10&30\\
\texttt{SL3.n1}&q\not\equiv1\pmod3&5&8&8&30\\
\texttt{SU3.2}&q\equiv2\pmod3&7&17&10&42\\
\texttt{SU3.n2}&q\not\equiv2\pmod3&7&8&8&42
\end{array}                                                     \tag{A2.2}
\]
The first column is the entire central class type.  In the tables
with center of order three it is parameterized by $a=0,1,2$;
its value need not equal the degree until $a=0$ is substituted.
The tables of types $17\times10$ are rectangular because a type
can parameterize multiple distinct characters or classes.

Every noncentral entry has the exact form
$z(q)=\sum_\alpha p_\alpha(q)u_\alpha$ with
$p_\alpha\in\mathbf Q[q]$ and $|u_\alpha|=1$ at every admissible
specialization of its integral parameters.  Set $x=q-q_0\ge0$ and
write
\[
 p_\alpha(q_0+x)=\sum_k c_{\alpha k}x^k,
 \qquad d(q_0+x)=\sum_k d_kx^k.
\]
Then
\[
 |z(q)|\le\sum_k\left(\sum_\alpha|c_{\alpha k}|\right)x^k.
                                                                  \tag{A2.3}
\]
For every nontrivial row and every column other than the first,
the exact coefficient comparison is
\[
       d_k/4-\sum_\alpha|c_{\alpha k}|\ge0
                         \quad\text{for every }k.          \tag{A2.4}
\]
All $386$ comparisons pass across all $46$ nontrivial row types.
Every translated degree polynomial has nonnegative coefficients,
with constant term at least the stated minimum.  Every class-length
polynomial except the first has nonnegative translated coefficients
and constant term greater than one.  Thus the omitted column is
exactly the center, and (A2.3)--(A2.4) prove (A2.1) for every
admissible field size and every class and character parameter.

The finite certificate is fully retained.  The four original files
have respective SHA-256 hashes
\begin{quote}\small\ttfamily
156d7173feac9283e0c3afb74c8c91c0bb13a1110bff274bc6d19be234b59c50\\
a5d5c0006bded69a9164a727104562226a6398dbe031bed15ba99e11fa58b7ec\\
364950438f2a8dac169b1834d43ec8957aae1781df1cfe9485834febf8846acc\\
f64356e81ec3f5fd951d551122cf0880dc7a8ebf352ec298148fe70695bce9b4.
\end{quote}
The retained original expressions, degrees and class lengths are in
\texttt{research/a2\_cores/generic\_inputs.json}, SHA-256
\texttt{4df14adc254b80d4d7a03ea18e9449c1e05c1ba72f2ad8b713067c8f252b01a2}.
The source extractor verifies these original hashes and executes
no downloaded Julia code.  Its identity check uses a restricted
arithmetic parser to verify that every central phase exponent
vanishes at $a=0$ and that the resulting entry is the separately
declared degree.  The previously reviewed formal-phase parser
collects exact rational polynomial coefficients.  The deterministic
program \texttt{build\_certificate.py} records every coefficient
in (A2.4), every degree polynomial and every translated degree.
Its output \texttt{generic\_certificate.json} has SHA-256
\texttt{7a8ec14dfbec9f2dbd29127c7aadbaa0bae46749dfa2375c1ba3a064318fa4db}.

The input boundary matters.  The corrected implementation includes
all central characters, not merely the third with trivial central
character in a misprinted header of Table~1b.  It uses the full
parameter moduli $q^2+q+1$ and $q^2-q+1$, the corrected phases on
the split cuspidal rows, and the corrected unitary class representative
whose eigenvalues have determinant one.  The linear mixed-torus
phase and the parameter exclusions are also corrected.  These
corrections are documented in the frozen source metadata; the
proof uses those complete corrected tables and does not certify an
uncorrected transcription of the printed tables.  The primary
character-table classification is a cited input, not a new result
of this certificate.
\end{proof}

\subsection*{D-A2-02. Every central cover and compatible normalizer}

\begin{lemma}[D-A2-02]
\label{D-A2-02}
Let $N$ act faithfully and irreducibly in degree $a$, where $N$ is
quasisimple with simple quotient $\operatorname{PSL}_3(q)$ or
$\operatorname{PSU}_3(q)$, $q\ge3$.  Then $a\ge6$, and every
nonscalar finite-order element of its complex linear normalizer
satisfies $\nu(T)\ge3a/16$.

If the simple quotient is one of
\[
 \operatorname{PSL}_3(3),\quad\operatorname{PSL}_3(4),\quad
 \operatorname{PSU}_3(3),\quad\operatorname{PSU}_3(4),\quad
 \operatorname{PSU}_3(5),                                  \tag{A2.5}
\]
the stronger estimate $\nu(T)\ge a/4$ holds.
For every other quotient the respective minimum $30$ or $42$
of D-A2-01 applies.
\end{lemma}

\begin{proof}
The complete exceptional-multiplier list
\cite[Chapter~5, Table~1, numbered p.~35]{BI:HHM} shows that the
only exceptional multiplier in the stated range is that of
$\operatorname{PSL}_3(4)$.  Otherwise the universal perfect cover
is the ordinary special linear or unitary group.  Every
representation of a central quotient lifts to such a cover.
For the generic ranges, D-A2-01 and the commutator argument
D-BI-18 give $\nu(T)\ge3a/16$: choose $u$ with $[T,u]\ne1$,
use $\operatorname{rank}([T,u]-1)\le2\nu(T)$, and compare
the trace of this commutator with $a/4$.  If it is a nonidentity
scalar, its rank is $a$, which gives a stronger inequality.
Thus this argument covers all compatible outer automorphisms
without any enumeration of generic outer cosets.

Here is the complete small-group calculation, using the official
CTblLib tables \cite{SPOR:CTblLib}.  Apart from $\operatorname{PSL}_3(4)$
the required tables and exact outcomes are
\[
\begin{array}{l|r|r|r|c}
\text{table}&\text{selected rows}&\text{nonscalar spectra}&
 \min a&\min\nu/a\\\hline
\texttt{L3(3)}&11&121&12&1/3\\
\texttt{L3(3).2}&10&140&12&1/3\\
\texttt{U3(3)}&13&169&6&2/7\\
\texttt{U3(3).2}&10&150&6&1/3\\
\texttt{U3(4)}&21&441&12&1/3\\
\texttt{U3(4).2}&10&190&12&1/3\\
\texttt{U3(4).4}&12&252&12&1/3\\
\texttt{3.U3(5)}&39&1443&20&8/21\\
\texttt{3.U3(5).2}&14&420&20&8/21\\
\texttt{3.U3(5).3}&90&8370&20&8/21
\end{array}                                                     \tag{A2.6}
\]
Their multiplier and outer-group data are respectively
$(1,C_2),(1,C_2),(1,C_4),(C_3,S_3)$, with the first pair for
$\operatorname{PSL}_3(3)$ and the others in increasing unitary
field size.  These are retained as the official extension metadata.
Every cyclic subgroup of each outer group is represented in (A2.6).
An invariant irreducible representation extends to any fixed
cyclic extension after rescaling its implementing matrix: if a
coset generator $b$ satisfies $b^d=u_0$, Schur's lemma says
$T^d=c\rho(u_0)$, and a scalar $d$th root corrects the relation.
Hence every compatible normalizer has, up to scalar, one of the
tested spectra.  This argument permits nonsplit cyclic extensions
and arbitrary central kernels.

For $S=\operatorname{PSL}_3(4)$ the universal multiplier is
$C_4^2\times C_3$ and $\operatorname{Out}(S)\simeq D_{12}$.
The new cyclic-extension and isoclinism-transfer lemmas
L34--1--3, proved in the adjoining section, reduce every actual
normalizing matrix to the complete published list of $36$
bicyclic extensions in \cite[\S3.3.1, Lemma~3.1]{D:BMO2017}.
Their core types are
\[
 S,\quad2.S,\quad3.S,\quad4_1.S,\quad4_2.S,\quad6.S,
 \quad12_1.S,\quad12_2.S.                                \tag{A2.7}
\]
For $S$ and $3.S$ the tables include the core, all three
involution extensions, and the order-three and order-six extensions.
For each other core they include the core and all three involution
extensions.  All $36$ tables are checked.  The smallest degree is
$6$, and every selected nonscalar spectrum satisfies the stronger
bound $\nu/a\ge3/10$.  Both order-four cover types and every
involution type are retained, so the source's discrepancy between
one displayed multiplier action and its following cover-label prose
does not affect coverage.

For completeness the computed universal-cover table
\texttt{(4\^{}2x3).L3(4)} is also included: its $303$ nonlinear
rows give $77568$ nonscalar spectra, minimum degree $6$ and
minimum ratio $1/3$.  Its retained provenance is construction by
\texttt{PossibleCharacterTablesOfTypeV4G}, rather than an ATLAS
source table.  It is a redundant consistency check; every essential
degree and support bound follows already from the eight cyclic-center
ATLAS core types and their published extension list.
No stored universal-cover-to-quotient fusion was available in GAP,
and none is claimed or used: completeness is supplied by L34--1--3
and the published $36$-extension classification.

The resulting complete input has $47$ tables, $1603$ selected
rows and $141024$ nonscalar spectra.  The exporter
\texttt{export\_small.g} retains every character, class size,
element order, prime power map, source string, and the actual
core-to-extension class fusions; it verifies full character-row
orthogonality.  The exact standalone verifier
\texttt{check\_small.py} selects every nonlinear core row and
every extension row whose restriction is an irreducible core row,
using exact equality of all character values.  Its checks include
degree-square and class-size sums, all $34$ core-to-extension
fusions, preservation of orders and power maps, trivial identity
kernel, normal class-size unions, and the full cyclic quotient
order detected by linear characters.  For an element of order $o$
it reconstructs all eigenvalue multiplicities by
\[
 m_b=\frac1o\sum_{j=0}^{o-1}\chi(g^j)\zeta_o^{-bj}
       \quad(0\le b<o),                                  \tag{A2.8}
\]
using exact cyclotomic-polynomial arithmetic.  Each multiplicity
is checked to be a nonnegative integer and their sum to be the
degree.  Scalar classes are exactly those with
$|\chi(g)|^2=\chi(1)^2$.  All the stated lower bounds pass.
The frozen \texttt{small\_inputs.json} has SHA-256
\texttt{3f356063efc3a3987126ddbbb6d12b0540fd71151f93eb4e83f4c3f6c26607e8}.
The individual row minima and spectra counts are retained in
\texttt{small\_checked.json}; the full row selection is not
replaced by a sample of low degrees.  This proves the lemma.
\end{proof}

\subsection*{D-A2-03. Exclusion of the entire supplier family}

\begin{theorem}[D-A2-03]
\label{D-A2-03}
Let an inducing block of a finite-cover Weil-sixfold source have
finite image $J\subset\operatorname{GL}(V)$.  Suppose it has an
irreducibly acting normal central product of quasisimple groups
whose simple quotients are $\operatorname{PSL}_3(q)$ or
$\operatorname{PSU}_3(q)$, with $q\ge3$ in every factor.
Then the actual period rank is at most eight.

Mixed linear and unitary types and field sizes, scalar phases,
compatible outer automorphisms and factor permutations are allowed.
No essential-self-duality, primitivity or tensor-indecomposability
assumption is needed, and there is no restriction on the inducing
character field or the polarization determinant.
\end{theorem}

\begin{proof}
Write $V=\bigotimes_{i=1}^t V_i$, with faithful irreducible
factor dimensions $a_i\ge6$, and $e=\prod_i a_i$.
The simple factors and their generated full matrix algebras are
intrinsic, so $J$ acts through factor intertwiners and compatible
factor permutations.  If a nonscalar element preserves every
factor, one intertwiner is nonscalar.  The largest eigenvalue
multiplicity fraction of a tensor product is at most that of
any one nonscalar factor.  D-A2-02 gives $\nu(g)\ge3e/16$.
For a nontrivial permutation cycle, the weighted cyclic-word
argument of D-SPOR-03, spelled out in D-G2T-03, gives
$\nu(g)\ge e/4$: after basis changes there is cyclic rotation
with one diagonal twist, and each word orbit contributes at
most one vector to a specified eigenspace.  Constant words
number $a$, while every other orbit has size at least two,
so the largest eigenspace is at most $(a^l+a)/2\le3a^l/4$.
Nonidentity scalar inertia has ordinary moving rank $e$.

If the actual period rank were nine, the intermediate-cover
factorization D-BI-01 requires genus zero and at least twelve
genuine branch points.  Indeed the multiplicity-six Euler formula
gives total moving rank $D=6+2e(1-h)$; for $h\ge2$ this is
negative since $e\ge6$, while for $h=1$ there are at most six
branch points and the intermediate Hurwitz image has dimension
at most six.  For $h=0$ the source dimension condition gives
$R\ge12$ and $D=2e+6$.  Therefore
\[
       12\left\lceil\frac{3e}{16}\right\rceil\le2e+6,
                     \qquad e\le24.                     \tag{A2.9}
\]
If $t\ge2$ then $e\ge36$, a contradiction.  For a single
factor in the generic range the lower bounds $30$ or $42$
also contradict (A2.9).  For one of the small quotients (A2.5),
the quarter-support bound improves the necessary inequality to
$12\lceil e/4\rceil\le2e+6$, hence $e\le6$.
The only remaining integer degree is $e=6$, for which the
left side is $24$ and the right side is $18$, again impossible.
This excludes every case and proves the actual-rank bound.
\end{proof}

\paragraph{Geometric and programme meaning.}
The theorem removes all these possible curve suppliers from the
remaining nonsplit-sixfold gate Q2i.  It gives a rank obstruction
for their actual variations of Hodge structure and associated
isotypic $\ell$-adic local systems, hence for this attempted
motivic supplier.  It produces no detecting cycle and no new
numerical bound by itself.  In the project's language it restricts
the supplier needed for the integral uniform form of condition~A
(cote~16, scan~050; SIGN\_NEXT SN-C03/C04); positivity, condition~I
(scan~049; cote~44, Th.~4), cannot repair the absent rank-nine
supplier.  The existing $K$-theory projectors and Cayley--Hamilton
division step of scans~048--049 still apply to sources that
actually supply cycles.  No conclusion about a ghost follows
from excluding this one source class.
